\maketitle
\thispagestyle{fancy}
\begin{abstract}
We compare processed gene expression, sparse dictionary codes, and frozen
scGPT embeddings for human pancreatic cell-type classification using the same
regularized multinomial logistic model family. On a fixed stratified split of
8,506 Baron cells, test Macro-F1 is 0.9939 for expression, 0.9810 for a
64-atom dictionary, and 0.9624 for a small custom scGPT checkpoint. A separate
retrospective nested donor-held-out analysis matches gene availability,
standardizes classifier inputs, and varies dictionary initialization. Its
respective donor-mean Macro-F1 scores are 0.9246, 0.9215, and 0.8823.
Dictionary predictions vary little across the tested initializations, whereas
matched atoms are only moderately reproducible. We relate these observations
to sparse-coding optimality, factor nonidentifiability, and the effect of
scaling on regularization. The results support sparse compression and the
importance of a strong expression baseline within this study; they do not
establish biological pathway recovery or general superiority over pretrained
single-cell models.
\end{abstract}

\section{Introduction and related work}
Single-cell RNA sequencing represents each cell through thousands of gene
measurements. A representation can be useful because it improves prediction,
compresses the data, or exposes patterns that can be inspected. These are
separate goals: an accurate classifier need not reveal stable biological
programs, and low reconstruction error need not imply useful cell-type
information. We ask whether a frozen pretrained representation provides more
linearly accessible cell-type information than a sparse representation learned
from the target training cells, and whether any compression is accompanied by
reproducible gene-level structure.

We compare three feature sets under the same logistic model family: processed
expression, signed sparse dictionary codes, and frozen scGPT embeddings. The
expression baseline tests whether representation learning is necessary for this
task. Our contributions are a provenance-checked comparison with frozen test
evaluation; a retrospective donor, scaling, and initialization analysis; and
an explicit connection between the optimization objectives, numerical
certificates, and the limits of interpretation. The mathematics is established
optimization theory applied to this experiment, not a claim of new theorems.

Sparse coding learns reusable components with sparse activities
\cite{olshausen1996}; fixing the dictionary yields an $\ell_1$-penalized
regression problem \cite{tibshirani1996}. Online dictionary learning provides
an established approach to the joint factorization \cite{mairal2010}, while
proximal methods explain the thresholding behavior of its convex coding
subproblem \cite{beck2009}. For single-cell data, consensus nonnegative matrix
factorization has been used to study identity and activity programs
\cite{kotliar2019}. Our signed atoms differ from nonnegative programs and
require sign/permutation alignment before stability comparisons.

scGPT motivates pretrained representations for single-cell tasks
\cite{cui2024}; work on pretraining size and diversity motivates examining
which benefits transfer to downstream evaluation \cite{denadel2026}. Our
comparison uses a small custom checkpoint trained by the project team,
not the original published large scGPT weights. Conclusions therefore concern
this checkpoint, this cohort, and regularized linear probing.
